\documentclass[11pt, a4paper]{article}

%% >>> preamble
\usepackage[T1]{fontenc}
\usepackage[a4paper, margin=2.5cm]{geometry}
\usepackage{amsmath}
\usepackage{graphicx}
\usepackage{booktabs}
\usepackage{siunitx}
\usepackage{fancyhdr}
\usepackage{xcolor}
\usepackage[colorlinks=true, allcolors=blue!60!black]{hyperref}
\pagestyle{fancy}
\fancyhf{}
\lhead{Lab Report -- Simple Pendulum}
\rhead{\today}
\cfoot{\thepage}
\setlength{\headheight}{14pt}
%% <<< preamble

\begin{document}
%% >>> snippet
\begin{titlepage}
  \centering
  \vspace*{3cm}
  {\scshape\LARGE University Name\par}
  \vspace{1cm}
  {\scshape\Large PHYS 101 -- Experimental Physics\par}
  \vspace{2.5cm}
  \rule{\linewidth}{0.6pt}\par\vspace{0.4cm}
  {\huge\bfseries Determining $g$ with a Simple Pendulum\par}
  \vspace{0.3cm}\rule{\linewidth}{0.6pt}\par
  \vspace{2cm}
  {\Large Your Name\par}
  {\large Student ID: 12345678\par}
  \vfill
  {\large
  Lab partner: A.~Partner\par
  Instructor: Dr.~Supervisor\par
  Date performed: \today\par}
\end{titlepage}

\section{Aim}
To measure the gravitational acceleration $g$ from the period of a simple pendulum.

\section{Theory}
For small oscillations the period is
\begin{equation}\label{eq:period}
  T = 2\pi \sqrt{\frac{L}{g}}
  \quad\Longrightarrow\quad
  g = \frac{4\pi^2 L}{T^2}.
\end{equation}

\section{Method}
Describe the apparatus and the procedure step by step.

\section{Results}
\begin{table}[htbp]
  \centering
  \caption{Measured periods for different lengths.}
  \label{tab:data}
  \begin{tabular}{@{}S[table-format=1.2] S[table-format=1.3] S[table-format=1.2]@{}}
    \toprule
    {$L$ (\unit{\metre})} & {$T$ (\unit{\second})} & {$g$ (\unit{\metre\per\second\squared})} \\
    \midrule
    0.50 & 1.419 & 9.80 \\
    0.75 & 1.738 & 9.80 \\
    1.00 & 2.006 & 9.81 \\
    \bottomrule
  \end{tabular}
\end{table}

\section{Discussion}
Compare with the accepted value $g = \qty{9.81}{\metre\per\second\squared}$ and
discuss sources of error.

\section{Conclusion}
State the final result with its uncertainty.
%% <<< snippet
\end{document}
